\cleardoublepage
\thispagestyle{plain}

% Inhaltsverzeichnis-Eintrag
\titlecontents{chapter}
[0em]
{\vspace{12pt}}
{\contentslabel{1em}}
{}
{\vzPunkte\contentspage}

% Code für Linksbuendigkeit der Klammern 
\makeatletter 
	\renewcommand*{\@biblabel}[1]{\makebox[\labelwidth][l]{[#1]}}
\makeatother

% Label-Width - Eintrag entspricht max. Anzahl an Ziffern
\setbiblabelwidth{99}

% Literaturverzeichnis Stil
\bibliographystyle{formatting/bibstyleFAPS}
%	Pfad:		C:\Program Files (x86)\MiKTeX 2.9\miktex\bin\bibtex.exe
%	Argumente:	"%bm"}
% Am besten aus Citavi exportieren
% \bibliography{Verzeichnis_Literatur}
% \bibliography{bibliography/literature/control,%
%              bibliography/literature/learning,%
%              bibliography/literature/robotics,%
%              bibliography/literature/simulation,%
%               bibliography/literature/cloth }
\bibliography{bibliography/literature}

%%%%%% Extra Verzeichnisse %%%%%
\cleardoublepage
\thispagestyle{plain}

% Inhaltsverzeichnis-Eintrag
\titlecontents{chapter}
[0em]
{}
{\contentslabel{1em}}
{}
{\vzPunkte\contentspage}


%%%%% Quellenverzeichnis %%%%%%

% Label-Width - Eintrag entspricht max. Anzahl an Ziffern
\setbiblabelwidth{9}

% Literaturverzeichnis Stil
\bibliographystyleQ{formatting/bibstyleFAPS}
% Am besten aus Citavi exportieren
% \bibliographyQ{Verzeichnis_Quellen}
\bibliographyQ{bibliography/sources}
%	Postprozessor:
		%Anwendung: 	bibtex.exe
		%Argumente:		Q



%% Layout Inhaltsverzeichnis wiederherstellen
%\input{Format_ToC}